% Options for packages loaded elsewhere
\PassOptionsToPackage{unicode}{hyperref}
\PassOptionsToPackage{hyphens}{url}
\PassOptionsToPackage{dvipsnames,svgnames,x11names}{xcolor}
%
\documentclass[
  11pt,
  a4paper,
]{ltjsarticle}
\usepackage{amsmath,amssymb}
\usepackage{iftex}
\ifPDFTeX
  \usepackage[T1]{fontenc}
  \usepackage[utf8]{inputenc}
  \usepackage{textcomp} % provide euro and other symbols
\else % if luatex or xetex
  \usepackage{unicode-math} % this also loads fontspec
  \defaultfontfeatures{Scale=MatchLowercase}
  \defaultfontfeatures[\rmfamily]{Ligatures=TeX,Scale=1}
\fi
\usepackage{lmodern}
\ifPDFTeX\else
  % xetex/luatex font selection
\fi
% Use upquote if available, for straight quotes in verbatim environments
\IfFileExists{upquote.sty}{\usepackage{upquote}}{}
\IfFileExists{microtype.sty}{% use microtype if available
  \usepackage[]{microtype}
  \UseMicrotypeSet[protrusion]{basicmath} % disable protrusion for tt fonts
}{}
\makeatletter
\@ifundefined{KOMAClassName}{% if non-KOMA class
  \IfFileExists{parskip.sty}{%
    \usepackage{parskip}
  }{% else
    \setlength{\parindent}{0pt}
    \setlength{\parskip}{6pt plus 2pt minus 1pt}}
}{% if KOMA class
  \KOMAoptions{parskip=half}}
\makeatother
\usepackage{xcolor}
\usepackage[margin=24mm]{geometry}
\setlength{\emergencystretch}{3em} % prevent overfull lines
\providecommand{\tightlist}{%
  \setlength{\itemsep}{0pt}\setlength{\parskip}{0pt}}
\setcounter{secnumdepth}{5}
\ifLuaTeX
  \usepackage{selnolig}  % disable illegal ligatures
\fi
\IfFileExists{bookmark.sty}{\usepackage{bookmark}}{\usepackage{hyperref}}
\IfFileExists{xurl.sty}{\usepackage{xurl}}{} % add URL line breaks if available
\urlstyle{same}
\hypersetup{
  colorlinks=true,
  linkcolor={blue},
  filecolor={Maroon},
  citecolor={Blue},
  urlcolor={blue},
  pdfcreator={LaTeX via pandoc}}

\author{}
\date{}

\begin{document}

\hypertarget{moorepenrose-ux64ecux4f3cux9006ux884cux5217}{%
\section{Moore--Penrose
擬似逆行列}\label{moorepenrose-ux64ecux4f3cux9006ux884cux5217}}

SVD

\[
A=U\Sigma V^\top
\]

に対し、非零特異値を逆数にした \texttt{Σ\^{}+} を作って

\[
\boxed{A^+=V\Sigma^+U^\top}
\]

と定義します。

\hypertarget{ux666eux901aux306eux9006ux884cux5217ux3068ux306eux95a2ux4fc2}{%
\subsection{普通の逆行列との関係}\label{ux666eux901aux306eux9006ux884cux5217ux3068ux306eux95a2ux4fc2}}

\texttt{A} が正方かつ可逆なら

\[
A^+=A^{-1}
\]

です。擬似逆は逆行列を長方形行列やランク落ち行列へ拡張した概念です。

\hypertarget{ux6700ux5c0fux4e8cux4e57ux89e3}{%
\subsection{最小二乗解}\label{ux6700ux5c0fux4e8cux4e57ux89e3}}

\[
\boxed{x_*=A^+b}
\]

は最小二乗解になります。解が複数ある場合には、その中で Euclid
ノルムが最小の解を選びます。

\hypertarget{svd-ux304bux3089ux898bux308bux610fux5473}{%
\subsection{SVD
から見る意味}\label{svd-ux304bux3089ux898bux308bux610fux5473}}

\texttt{U\^{}Tb} で出力を特異方向へ分解し、各非零成分を \texttt{1/σ\_i}
倍して逆変換し、\texttt{V} で入力空間へ戻します。

\texttt{σ\_i=0}
の方向は元の変換で完全に失われているため、逆算できません。擬似逆はその方向を0として扱います。

\end{document}
